\documentclass[12pt, a4]{article}

\usepackage{amsmath}
\usepackage{german}
\usepackage{url}
\usepackage[dvips]{epsfig}

%\graphicspath{{/home/ritz/Teaching/fcn/fig/}}

%----------------- layout ----------------------------------------

\oddsidemargin=0in
\textwidth=6in
\topmargin=-0.7in
\textheight=9.5in
\baselineskip=16pt
\font\mittel=cmr17
\font\roman=cmr12
\font\klein=cmr9
\font\mgross=cmr10
\newfont{\CAP}{cmcsc10 scaled 900}
\pagestyle{empty}
\parindent=0mm

%----------------- macros ----------------------------------------
\newcommand{\conc}[1]{[{\rm C}]_{\text{#1}}}
\newcommand{\units}[1]{\text{#1}}
\newcommand{\kalium}{$\text{K}^+$}
\newcommand{\natrium}{$\text{Na}^+$}
\newcommand{\chlor}{$\text{Cl}^-$}
\newcommand{\wasser}{$\text{H}_2\text{O}$}
%-----------------------------------------------------------------


\begin{document}


\parbox{2cm}{
\epsfig{file = bccnlogo.eps, width = 1.8cm}
}
\parbox{11cm}{
\begin{center}
\large HUMBOLDT-UNIVERSIT"AT \hskip 0.1 cm ZU \hskip 0.1 cm BERLIN
\vskip 0.1 true cm
\mgross BERNSTEIN CENTRE FOR COMPUTATIONAL NEUROSCIENCE
\end{center}
}
\parbox{2cm}
{
\hfill
\epsfig{file = hublogo.eps, width = 1.8cm}
}

\vskip 0.6 true cm

\leftline{\CAP Humboldt-Universit"at zu Berlin 
\hfill Phone: 030/2093-9110}
\leftline{\CAP Philippstr. 13 House 6
              \hfill Fax: 030/2093-6771}
\leftline{\CAP \hfill webpage: http://www.bccn-berlin.de/}




\vskip 0.8 true cm
\centerline{\bf Models of Neural Systems I, WS 2010/11}
\centerline{\bf Computer Practical 1}


\vskip 0.8 true cm
{\bf First steps in Python}

\medskip
Python is a modern high-level programming language. In contrast to
C/C++ and Java, Python is an interpreted language, which means that
Python programs (also called scripts) don't have to be compiled. It
makes programming Python very quick and interactive. It also allows
you to see immediately the results of all operations, analyse them or
further process them. Here you will learn the basic concepts of
Python.\\

\begin{enumerate}
\item  Open an interactive shell (we recommend to use
    \texttt{ipython}, see below for a short introduction). Implement  a simple (!) calculator:
\begin{verbatim}
a = 3
b = 4
print 'a+b=',a+b
\end{verbatim}
Open your favourite editor, paste the code and save it as
\texttt{calc.py}. Run the script within the \texttt{ipython}
shell (using \texttt{\%run} command, see also below).\\

%Python shell, scripts
\item \textbf{Data structures}.
    
    \begin{enumerate}
        \item Create variables containing: a number,
    a string and a list. 
        \item Find out what methods are available for strings and
            lists (use python \texttt{dir} function).
            \textit{Hint:}
    IPython offers convenient tab-completion which can be used to look
    up methods names.
         \item Try adding two lists and two strings together. What is
             the resulting data?
\end{enumerate}
%data structures, classes, introspection
\item \textbf{Functions}. 
    
\begin{enumerate} \item  Write a function
    which returns first $N$ elements of Fibbonacci series that
    is defined by the following recursive formula:
            
\begin{equation}
F(n)=F(n-1)+F(n-2),\rm{ for}\, n\geq2. 
\end{equation}
The first elements of the series are $F(0)=0$ and
$F(1)=1$.

To solve the problem, you can either use recursion (i.e. a function calling
itself) or a for-loop.

\item \textit{Documenting your code.} Write a documentation for your
    function.
    
\item \textit{Testing your code.} Test the function with several
    predefined arguments and compares its result to a hard-coded value
    (e.g. F(0)=0, F(1)=1, F(3)=3, etc.).
% defining functions, indentation, flow control (while loop), doc
% string,  recurssion
\end{enumerate}

\item \textbf{For loop and list comprehensions.} 

        Generate following sequences:

        \begin{itemize}
            \item odd numbers from 1 to 11,
            \item letters from ''a'' to ''k'',
            \item 10 first powers of 2 (i.e. 2, 4,
                8\dots),
        \end{itemize}


\end{enumerate}

\vskip 0.8 true cm
{\bf IPython}

\medskip
In order to make programming in Python more efficient we recommend you
to use \texttt{ipython}. IPython is an extended interactive shell (a
replacement for \texttt{python}) which supports advanced line editing
(history, tab completion) and special commands (magic commands). The
list of the most commonly used magic commands is below: \\

\texttt{\%magic} -- a full list of magic functions, \\
\texttt{\%hist} -- show input history, \\
\texttt{\%run} -- run an external python script inside IPython, \\
\texttt{\%edit} -- open an external editor (default \texttt{vim}), \\
\texttt{\%pdb} -- run a debugger when an error occurs, \\
\texttt{\%cpaste} -- paste Python code (for example from clipboard),\\
\texttt{cmd?} -- print a documentation of a \texttt{cmd}, \\
\texttt{cmd??} -- print a documentation of a \texttt{cmd} together
with the source code.\\

\medskip
{\bf Links}
\begin{itemize}
    \item Python Tutorial: \url{http://docs.python.org/tutorial/}
    \item Python Scientific: 
        \url{https://portal.g-node.org/python-autumnschool/_media/pythonscientific.pdf}
\end{itemize}


\vfill
\centerline{\CAP Contact}
\CAP

Urs Bergmann \hfill ubergmann@fias.uni-frankfurt.de \\ 
Matthias Guggenmos\hfill Email: matthias.guggenmos@bccn-berlin.de\\
\underline{Richard Kempter} \hfill Phone: 2093-8925 \hfill
Email: r.kempter(at)biologie.hu-berlin.de \\
Robert Schmidt \hfill Phone: 2093-8926 \hfill
Email: r.schmidt(at)biologie.hu-berlin.de \\
Bartosz Telenczuk\hfill Phone: 2093-8838 \hfill
Email: b.telenczuk(at)biologie.hu-berlin.de \\
\end{document}
 
